% d3 — Nguyên lý sensing element: mặt cắt ngang (vách cứng + sandwich + ngón ép + kéo trước Fp)
\documentclass[border=4pt]{standalone}
\usepackage{fontspec}
\setmainfont{Liberation Serif}
\usepackage{tikz}
\usetikzlibrary{patterns}
\begin{document}
\begin{tikzpicture}[font=\footnotesize]
% ---- vách khung cứng (in 3D): nền cố định, ký hiệu gạch chéo
\fill[pattern=north east lines, pattern color=black!35] (-1.6,-0.42) rectangle (1.6,0.0);
\draw[black!60] (-1.6,-0.42) rectangle (1.6,0.0);
% ---- sandwich: điện cực đồng (dưới) - Velostat (giữa) - điện cực đồng (trên)
\fill[orange!75] (-1.6,0.00) rectangle (1.6,0.20);
\fill[gray!22]  (-1.6,0.20) rectangle (1.6,0.80);
\fill[orange!75] (-1.6,0.80) rectangle (1.6,1.00);
% ---- vùng nén A_c: dải highlight TRÊN HÀNG MÀNG (0.20..0.80), đúng chỗ ngón đè
\fill[orange!30, opacity=0.6] (-0.9,0.20) rectangle (0.9,0.80);
% ---- nắp kẹp kéo trước (gạch chéo khác hướng)
\fill[pattern=north west lines, pattern color=black!30] (-1.6,1.00) rectangle (1.6,1.18);
\draw[black!60] (-1.6,1.00) rectangle (1.6,1.18);
% ---- hai vít xiết tạo lực kéo trước Fp
\fill[gray!45] (-1.34,1.18) rectangle (-1.14,1.52);
\fill[gray!45] ( 1.14,1.18) rectangle ( 1.34,1.52);
\draw[->,>=stealth,thick,black!70] (-1.24,2.10) -- (-1.24,1.56);
\draw[->,>=stealth,thick,black!70] ( 1.24,2.10) -- ( 1.24,1.56);
\node at (-1.24,2.32) {$F_p$};
\node at ( 1.24,2.32) {$F_p$};
% ---- đốt ngón ép
\draw[rounded corners=3pt, fill=blue!10, draw=blue!55!black] (-0.9,1.85) rectangle (0.9,2.75);
\node at (0,2.30) {đốt ngón};
% mặt ép = đáy ngón (nét đứt) + lực F đặt lệch phải (tránh thước t0 bên trái)
\draw[black!55,dashed] (-0.9,1.85) -- (0.9,1.85);
\draw[->,>=stealth,very thick,black!70] (0.45,1.84) -- (0.45,1.20);
\node[right] at (0.58,1.50) {$F$};

% ---- t0: bề dày màng, đo ở phần TRÁI màng (chỗ không bị nén, không đè highlight)
\draw[<->,>=stealth,black!60] (-1.35,0.22) -- (-1.35,0.78);
\node[font=\scriptsize,black!70] at (-1.08,0.50) {$t_0$};

% ---- nhãn trái (leader chạm đúng lớp)
\node[anchor=east,font=\scriptsize,black!75] at (-1.90,-0.10) {vách khung cứng (in 3D)};
\draw[black!55] (-1.84,-0.10) -- (-1.62,-0.10);
\node[anchor=east,font=\scriptsize,black!75] at (-1.90,0.10) {điện cực lá đồng (dưới)};
\draw[black!55] (-1.84,0.10) -- (-1.62,0.10);
\node[anchor=east,font=\scriptsize,black!75] at (-1.90,0.50) {vùng nén $A_c$};
\draw[black!55] (-1.84,0.50) -- (-0.90,0.50);
\node[anchor=east,font=\scriptsize,black!75] at (-1.90,1.52) {khe hở danh định};
\draw[black!55] (-1.84,1.52) -- (-0.40,1.52);
\node[anchor=east,font=\scriptsize,black!75] at (-1.90,1.86) {mặt ép của ngón};
\draw[black!55] (-1.84,1.86) -- (-0.92,1.86);
% ---- nhãn phải (anchor west, x >= 1.6 để không phủ mép stack)
\node[anchor=west,font=\scriptsize,black!75] at (1.75,0.66) {màng Velostat};
\draw[black!55] (1.34,0.66) -- (1.71,0.66);
\node[anchor=west,font=\scriptsize,black!75] at (1.75,0.90) {điện cực lá đồng (trên)};
\draw[black!55] (1.62,0.90) -- (1.71,0.90);
\node[anchor=west,font=\scriptsize,black!75] at (1.75,1.32) {nắp kẹp + vít kéo trước};
\draw[black!55] (1.30,1.32) -- (1.71,1.32);
\node[anchor=west,font=\scriptsize,black!75] at (1.75,-0.20) {khi bị nén: $R$ giảm ($\Delta R<0$)};
\draw[black!55] (1.62,-0.16) -- (1.71,-0.16);

% ---- chú giải cuối hình
\node[anchor=north west,font=\scriptsize,black!65,align=left] at (-2.60,-0.82)
  {$t_0$ -- bề dày màng \quad $A_c$ -- diện tích vùng nén\\
   $F_p$ -- lực kéo trước do hai vít xiết nắp kẹp \quad $F$ -- lực ngón ép vào mặt ép};
\end{tikzpicture}
\end{document}
